\section{Compensation over other class-number-one fields}

\begin{proposition}[HD-COMP: an effective contraction from a split family]
\label{HD-COMP}
Let $K$ be imaginary quadratic with class number one, write
$\mathcal O_K=\mathbf Z[\tau]$, and let $(A,L)$ have dimension $2n$,
signature $(n,n)$, and integral compatible $\mathcal O_K$ action.
Put $\Delta=\chi(L)>0$ and $c_K=1+\operatorname{Nm}(\tau)$.
Suppose that a bound $D_{n+1}(x)$ is proved for every split polarized
Weil $(2n+2)$-fold over $K$ with polarization product $x$.
Then, with no determinant restriction on $A$,
\[
 \delta_W(A,L)\le c_K D_{n+1}(\Delta^2).
\]
All cycles in this assertion are effective integral cycles. If the
higher-dimensional bound is numerical and explicit, so is this bound.
If it is only existential, this conclusion is only existential.
\end{proposition}
\begin{proof}
The normalized Hermitian form is
$h(x,y)=E(x,\tau y)-\bar\tau E(x,y)$, where $E$ is the alternating
polarization form. In an integral $\mathcal O_K$ basis its determinant
has absolute value $\Delta$ and sign $(-1)^n$.
Here is the normalization check for a general quadratic order basis.
Writing $\tau+\bar\tau=b$ and $\tau\bar\tau=c$, the real determinant
of a Hermitian matrix is the norm of its $K$-determinant. The map from
$h$ to $E$ is the coefficient of $\tau$ in $h$ (with the fixed choice
of conjugate-linear argument), and on a rank-one block $h=(a)$ the
integral basis $(v,\tau v)$ gives the alternating matrix
$\bigl(\begin{smallmatrix}0&a\\-a&0\end{smallmatrix}\bigr)$.
Change of $K$-basis changes both determinant formulas by the same
norm square. Thus $\det_{\mathbf Z}E=(\det_K h)^2$, without an
extra discriminant factor. The class-number-one assumption supplies
the free integral $\mathcal O_K$ basis.

Let $E_K=\mathbf C/\mathcal O_K$, with principal class $\vartheta$,
and put $T=E_K^2$ with balanced action and
$\lambda=\vartheta_1+\Delta\vartheta_2$. Its Hermitian determinant
is $-\Delta$. Consequently $(A\times T,H)$, where
$H=p_A^*c_1(L)+p_T^*\lambda$, has polarization product $\Delta^2$
and determinant $(-1)^{n+1}\Delta^2$. Rational Hermitian
classification, Schoen (1998), Section~7, printed p.~332, makes this
space split: its signed determinant is the norm of the rational
integer $\Delta$. It therefore has an effective detector $Z$ of
dimension $n+1$ and degree at most $D_{n+1}(\Delta^2)$.

On $T$, the two coordinate elliptic divisors and the graphs of
$1,\tau$ span the rational Neron--Severi group. Under the usual
Hermitian-matrix description, their off-diagonal entries span $K$
over $\mathbf Q$. Subtracting the corresponding diagonal classes
from the two graphs gives a basis of $W_K(T)$ for the balanced
action. For the graph of $\tau$ the diagonal coefficients are
$\operatorname{Nm}(\tau)$ and $1$, in an order determined by the
choice of coordinates. Its matrix has eigenvalues $0,c_K$; the
unit graph has eigenvalues $0,2$. Because $c_K\ge2$, every one of
these four effective divisors $D$ satisfies
\[
 [D]\le c_K(\vartheta_1+\vartheta_2)\le c_K\lambda
\]
in the nef order.

The nonzero Weil projection of $[Z]$ lies in the Kunneth component
$W_K(A)\otimes W_K(T)$. The intersection pairing on $W_K(T)$ is
nondegenerate. If contraction against each of the four effective
divisors were zero after projection to $W_K(A)$, contraction
against their two off-diagonal linear combinations would also be
zero, contradicting that nonzero tensor. Choose a divisor with
nonzero contraction, translate it generally to meet $Z$ properly,
and push $Z\cdot p_T^*D$ to $A$. This is an effective integral
$n$-cycle and detects $W_K(A)$. The incidence and translation
argument proving properness is exactly the one given in HD-E6;
it applies to every dimension. Its degree is at most
\[
 \int_Z p_A^*c_1(L)^n p_T^*[D]
 \le c_K\int_Z p_A^*c_1(L)^n p_T^*\lambda
 \le c_K\int_Z H^{n+1}
 \le c_K D_{n+1}(\Delta^2).
\]
Choosing a detecting reduced component only decreases this degree.
This proves the assertion without dividing a cycle or supposing
that $Z$ itself is a pure Weil cycle.
\end{proof}

\begin{corollary}[HD-COMP4: numerical fourfold bounds beyond the principal case]
\label{HD-COMP4}
For compatible integral maximal-order actions the following hold on
every fixed fourfold polarization datum, without a determinant
restriction:
\[
\begin{array}{c|c}
 K & \text{upper bound for }\delta_W(A,L)\\ \hline
 \mathbf Q(i) & 128\Delta^{12}D_{G,6}\\
 \mathbf Q(\sqrt{-7}) & 3D_{6,7}(\Delta^2)\\
 \mathbf Q(\sqrt{-2}) & 3D_{6,2}(\Delta^2).
\end{array}
\]
Here $D_{G,6}$ is the integer of SF-03, and $D_{6,7},D_{6,2}$
are the fully explicit functions of D-G24-05 and D-GF4-04.
For an integral action only of $\mathbf Z+f\mathcal O_K$, replace
$\Delta$ in these bounds by $f^8\Delta$. Thus the Gaussian
bound, for example, becomes $128f^{96}\Delta^{12}D_{G,6}$.
The much smaller principal Gaussian bound HD-G06 is retained
whenever its hypotheses hold.
\end{corollary}
\begin{proof}
For $K=\mathbf Q(i)$ choose $\tau=i$, so $c_K=2$ and apply
SF-04a to the split sixfold of polarization product $\Delta^2$:
$2(2\Delta^2)^6D_{G,6}=128\Delta^{12}D_{G,6}$.
For $K=\mathbf Q(\sqrt{-7})$ use an integral generator of norm
$2$, and for $K=\mathbf Q(\sqrt{-2})$ use $\tau=\sqrt{-2}$.
Both have $c_K=3$, so the two reconstructed sixfold bounds give
the remaining lines. For conductor $f$, put
$\Lambda'=f\mathcal O_K\Lambda$. Then
$f\Lambda\subset\Lambda'\subset\Lambda$ and
$I=[\Lambda:\Lambda']\le f^8$. The isogeny
$A'=V/\Lambda'\to A=V/\Lambda$ has pulled-back polarization
product $I\Delta$. Pushing its detector forward has exactly the
same degree relative to these polarizations, and keeps its nonzero
Weil projection. The displayed functions are increasing, so
$I\Delta\le f^8\Delta$ proves the stated bounds.
\end{proof}

The same contraction applied to D-SR-05 gives bounded division on
every fourfold determinant for each class-number-one field, using
the supplied Perry deformation input and without a Markman theorem
as a premise. At this stage its detector integer is existential for
the fields without a numerical split-sixfold bound. This is not
advertised as a numerical solution of those remaining gates.
In the project's frame these are new suppliers for the integral
condition A of cote~16 scan~050; their effective intersection uses
condition I of scan~049 and cote~44 Theorem~4 only after a higher
detector has been supplied. They are correspondences between the
weight-$2n+2$ and weight-$2n$ determinant realizations, with Tate
twists $(n+1)$ and $(n)$. The spread, $K_0$, and cycle-space
interpretations are those specified in the framework section.
